\documentclass[11pt]{article}\usepackage{layout}\begin{document}
\begin{canvas}\Header{第三章 · 全球开源影响力排行榜}{violet}\Footer{01}
\Box{15}{27}{180}{\fontsize{20}{27}\selectfont\color{ink}\sffamily \bfseries 第三章　全球开源影响力排行榜}
\Box{15}{42}{180}{\fontsize{14}{20}\selectfont\color{violet}\sffamily \bfseries 从主体位置走向持续协作能力}
\Box{15}{55}{180}{\fontsize{10.3}{16}\selectfont\color{ink}第二章回答开放生态在哪里生长，本章进一步识别哪些地区、组织、项目与模型形成了影响。排行榜提供寻找关键主体的入口；将位次与技术分工、实际接口和维护记录相连，才能理解开放资源如何转化为可持续公共能力。}
\Box{15}{84}{180}{\fontsize{11.3}{16}\selectfont\color{violet}\sffamily \bfseries 六类对象，一条递进的观察路径}
\Rule{15}{115}{195}{line}
\Box{15}{98}{13}{\fontsize{11}{15}\selectfont\color{violet}\sffamily \bfseries 3.1}
\Box{31}{98}{35}{\fontsize{11}{15}\selectfont\color{ink}\sffamily \bfseries 行政区划}
\Box{70}{98}{106}{\fontsize{10.2}{15}\selectfont\color{ink}从全球分布落到省州与城市}
\Box{70}{105}{106}{\fontsize{8.7}{12}\selectfont\color{muted}\sffamily 参与规模、影响位置与区域内的多城结构}
\Box{178}{98}{17}{\fontsize{8.8}{12}\selectfont\color{muted}\sffamily 02—04}
\Rule{15}{135}{195}{line}
\Box{15}{118}{13}{\fontsize{11}{15}\selectfont\color{violet}\sffamily \bfseries 3.2}
\Box{31}{118}{35}{\fontsize{11}{15}\selectfont\color{ink}\sffamily \bfseries 企业}
\Box{70}{118}{106}{\fontsize{10.2}{15}\selectfont\color{ink}辨认资源供给与持续投入的主体}
\Box{70}{125}{106}{\fontsize{8.7}{12}\selectfont\color{muted}\sffamily 全球布局、时序变化与中国企业分工}
\Box{178}{118}{17}{\fontsize{8.8}{12}\selectfont\color{muted}\sffamily 05—08}
\Rule{15}{155}{195}{line}
\Box{15}{138}{13}{\fontsize{11}{15}\selectfont\color{violet}\sffamily \bfseries 3.3}
\Box{31}{138}{35}{\fontsize{11}{15}\selectfont\color{ink}\sffamily \bfseries 基金会}
\Box{70}{138}{106}{\fontsize{10.2}{15}\selectfont\color{ink}识别跨组织共同建设的载体}
\Box{70}{145}{106}{\fontsize{8.7}{12}\selectfont\color{muted}\sffamily 项目组合、维护重心与智能体公共组件}
\Box{178}{138}{17}{\fontsize{8.8}{12}\selectfont\color{muted}\sffamily 09—11}
\Rule{15}{175}{195}{line}
\Box{15}{158}{13}{\fontsize{11}{15}\selectfont\color{violet}\sffamily \bfseries 3.4}
\Box{31}{158}{35}{\fontsize{11}{15}\selectfont\color{ink}\sffamily \bfseries 项目}
\Box{70}{158}{106}{\fontsize{10.2}{15}\selectfont\color{ink}从影响位次进入技术与协作关系}
\Box{70}{165}{106}{\fontsize{8.7}{12}\selectfont\color{muted}\sffamily 全球与中国结构、人员接点和兼容接口}
\Box{178}{158}{17}{\fontsize{8.8}{12}\selectfont\color{muted}\sffamily 12—15}
\Rule{15}{195}{195}{line}
\Box{15}{178}{13}{\fontsize{11}{15}\selectfont\color{violet}\sffamily \bfseries 3.5}
\Box{31}{178}{35}{\fontsize{11}{15}\selectfont\color{ink}\sffamily \bfseries 新势力}
\Box{70}{178}{106}{\fontsize{10.2}{15}\selectfont\color{ink}观察新增供给怎样形成体系}
\Box{70}{185}{106}{\fontsize{8.7}{12}\selectfont\color{muted}\sffamily 仓库与项目的不同结构、CANN 多仓协同}
\Box{178}{178}{17}{\fontsize{8.8}{12}\selectfont\color{muted}\sffamily 16—18}
\Rule{15}{215}{195}{line}
\Box{15}{198}{13}{\fontsize{11}{15}\selectfont\color{violet}\sffamily \bfseries 3.6}
\Box{31}{198}{35}{\fontsize{11}{15}\selectfont\color{ink}\sffamily \bfseries 模型}
\Box{70}{198}{106}{\fontsize{10.2}{15}\selectfont\color{ink}把代码协作扩展到模型接续}
\Box{70}{205}{106}{\fontsize{8.7}{12}\selectfont\color{muted}\sffamily 全球百强、中国供给及 Qwen / DeepSeek}
\Box{178}{198}{17}{\fontsize{8.8}{12}\selectfont\color{muted}\sffamily 19—21}
\Box{15}{225}{180}{\fontsize{11.3}{16}\selectfont\color{violet}\sffamily \bfseries 两类指标，分别解释网络中的位置}
\Box{15}{237}{86}{\fontsize{9.7}{15}\selectfont\color{ink}OpenRank：依据协作连接及相连节点的重要性衡量影响。人数、仓库数与分值分别阅读；跨组织标签可能重叠，不合成为一个总量。}
\Box{109}{237}{86}{\fontsize{9.7}{15}\selectfont\color{ink}ModelRank：依据模型派生与复用关系观察影响。本章使用 2026 年 6 月月榜，不与代码侧 OpenRank 或全年累计直接比较。}
\Box{15}{264}{180}{\fontsize{8.9}{13.2}\selectfont\color{muted}\sffamily 主统计期：2025-07-01—2026-06-30。原榜待核窗口、历史记录与期后观察各自注明。榜单呈现网络位置，不替代技术质量、开放许可或公共成效评价。}
\end{canvas}
\begin{canvas}\Header{第三章 · 全球开源影响力排行榜}{violet}\Footer{02}
\Box{15}{27}{180}{\fontsize{20}{27}\selectfont\color{ink}\sffamily \bfseries 3.1 全球行政区划影响力排行榜}
\Box{15}{40}{180}{\fontsize{9.7}{13.5}\selectfont\color{muted}\sffamily 第二章呈现生态的地域分布，本节进一步识别州、省与城市中的影响节点。地区排名连接参与基础与协作网络，为后续寻找企业、基金会和项目提供空间入口。}
\Box{15}{53}{180}{\includegraphics[width=180mm]{.build/assets/fig-3-01.pdf}}
\Box{15}{199}{180}{\fontsize{7.7}{10.7}\selectfont\color{muted}\sffamily 图 3-1　全球行政区影响力前 30 位、前 100 位国别覆盖及席位结构。来源：GOSD 2026 行政区划榜与国别汇总；原榜统计起止待核。州、省、构成国等单元不完全等价；席位指入榜地区数量。}
\Box{15}{214}{87}{\fontsize{10.1}{14}\selectfont\color{violet}\sffamily \bfseries ① 全球协作格局由具体地区共同构成}
\Box{15}{220}{87}{\fontsize{9.8}{14.9}\selectfont\color{ink}加利福尼亚州、英格兰和纽约州列前三，前十位来自七个国家。圣保罗州、安大略省、卡纳塔克邦、柏林与首尔也居前列：观察全球协作，需要同时看国家背景与具体地区承载的组织、项目。}
\Box{108}{214}{87}{\fontsize{10.1}{14}\selectfont\color{violet}\sffamily \bfseries ② 全球南方的参与具有区域落点}
\Box{108}{220}{87}{\fontsize{9.8}{14.9}\selectfont\color{ink}前 100 位覆盖 32 国，印度与巴西分别有 12 席、9 席；尼日利亚、肯尼亚与埃及各有 1 席。国际合作可据此继续寻找当地社区和维护者，将对全球南方潜力的判断落实到可接续的合作对象。}
\Box{15}{248}{87}{\fontsize{10.1}{14}\selectfont\color{violet}\sffamily \bfseries ③ 头部集聚与多地参与同时存在}
\Box{15}{254}{87}{\fontsize{9.8}{14.9}\selectfont\color{ink}美国、印度、巴西、德国和中国共占 61 席，其余 27 国合计 39 席。这表明高位节点既集中于少数国家，也延伸至更多地区；席位统计描述入榜分布，不能替代各国贡献或影响力份额。}
\Box{108}{248}{87}{\fontsize{10.1}{14}\selectfont\color{violet}\sffamily \bfseries ④ 节点数量与影响位置需要联读}
\Box{108}{254}{87}{\fontsize{9.8}{14.9}\selectfont\color{ink}美国在前 30 位有 11 席；英国仅英格兰一席，却居第 2。中国与德国各有 4 席，位次分布不同。数量说明节点覆盖，位次提示协作位置；下一页进一步把这种差异与参与规模放在同一组地区中比较。}
\end{canvas}
\begin{canvas}\Header{第三章 · 全球开源影响力排行榜}{violet}\Footer{03}
\Box{15}{27}{180}{\fontsize{20}{27}\selectfont\color{ink}\sffamily \bfseries 3.1.1 参与规模与影响位置}
\Box{15}{40}{180}{\fontsize{9.7}{13.5}\selectfont\color{muted}\sffamily 参与人数说明基础有多大，网络影响说明在协作中处于什么位置。同一组 30 个地区的两种排序显示，规模优势与影响位置并不重合，地方发展因此需要多维对标。}
\Box{15}{53}{180}{\includegraphics[width=180mm]{.build/assets/fig-3-02.pdf}}
\Box{15}{199}{180}{\fontsize{7.7}{10.7}\selectfont\color{muted}\sffamily 图 3-2　同一组 30 个地区的参与规模与影响位次。来源：GOSD 2026 行政区划榜；人数单位与覆盖范围待核，因此仅比较榜内排序。圆内数字对应图 3-1 名次；连线不表示时间变化。}
\Box{15}{214}{87}{\fontsize{10.1}{14}\selectfont\color{violet}\sffamily \bfseries ① 柏林显示规模与影响并不同序}
\Box{15}{220}{87}{\fontsize{9.8}{14.9}\selectfont\color{ink}柏林是图中的第 9 号点，参与规模列第 17、影响列第 9。圆内数字与前页名次相连；越靠左表示规模位次越高，越靠上表示影响位次越高。这种对照揭示不同观察维度，不是人均效率评价。}
\Box{108}{214}{87}{\fontsize{10.1}{14}\selectfont\color{violet}\sffamily \bfseries ② 两项前十交集为六地，另外四地不同}
\Box{108}{220}{87}{\fontsize{9.8}{14.9}\selectfont\color{ink}北京、上海、东京与马哈拉施特拉邦仅进入规模前十；华盛顿州、安大略省、柏林与首尔仅进入影响前十。人数与影响联读，能够同时识别参与基础较大的地区，以及协作影响位置较高的节点。}
\Box{15}{248}{87}{\fontsize{10.1}{14}\selectfont\color{violet}\sffamily \bfseries ③ 从位置差异走向项目与维护角色}
\Box{15}{254}{87}{\fontsize{9.8}{14.9}\selectfont\color{ink}华盛顿州的规模、影响位次分别为 13、4，安大略省为 14、6。下一步应追踪当地组织关联的关键项目、评审与维护角色、跨项目参与。地区排序提供调查线索，项目证据才能解释差异如何形成。}
\Box{108}{248}{87}{\fontsize{10.1}{14}\selectfont\color{violet}\sffamily \bfseries ④ 中国地区要把参与基础接入协作网络}
\Box{108}{254}{87}{\fontsize{9.8}{14.9}\selectfont\color{ink}北京、上海的规模位次为 5、9，影响位次为 14、16；广东为 11、25，台湾为 16、26。地方行动需要进一步进入具体组织与项目，支持评审、适配和持续维护，观察参与基础怎样获得更稳定的协作位置。}
\end{canvas}
\begin{canvas}\Header{第三章 · 全球开源影响力排行榜}{violet}\Footer{04}
\Box{15}{27}{180}{\fontsize{20}{27}\selectfont\color{ink}\sffamily \bfseries 3.1.2 中国省市：高位节点与多城结构}
\Box{15}{40}{180}{\fontsize{9.7}{13.5}\selectfont\color{muted}\sffamily 从全球节点转向中国，省级榜呈现区域整体，城市榜识别具体中心。两级观察共同揭示区域内部的层次，为连接当地企业、维护团队和跨城共同任务提供依据。}
\Box{15}{53}{180}{\includegraphics[width=180mm]{.build/assets/fig-3-03.pdf}}
\Box{15}{199}{180}{\fontsize{7.7}{10.7}\selectfont\color{muted}\sffamily 图 3-3　中国省市影响力、两级名次及分值距离。来源：GOSD 2026 省级榜与城市榜；原榜统计起止待核。③选取城市前 30 位中有多城入榜的七省；连线表示行政归属，省市分值不相加。}
\Box{15}{214}{87}{\fontsize{10.1}{14}\selectfont\color{violet}\sffamily \bfseries ① 省级梯次需要结合分值理解}
\Box{15}{220}{87}{\fontsize{9.8}{14.9}\selectfont\color{ink}北京、上海居前两位，广东、台湾和浙江随后；浙江分值约为江苏的 2.06 倍。相邻省级名次可能对应明显不同的分值规模，选择对标地区时，应继续比较产业、组织与项目构成。}
\Box{108}{214}{87}{\fontsize{10.1}{14}\selectfont\color{violet}\sffamily \bfseries ② 城市榜定位具体资源的汇聚点}
\Box{108}{220}{87}{\fontsize{9.8}{14.9}\selectfont\color{ink}北京、上海、杭州、深圳和台北位列城市榜前五。杭州城市排名第 3，浙江省级排名第 5；二者反映不同层级。后续企业与模型观察将继续辨认杭州的组织节点，省市排名本身不证明其 AI 产业优势成因。}
\Box{15}{248}{87}{\fontsize{10.1}{14}\selectfont\color{violet}\sffamily \bfseries ③ 多城结构各具层次，支持差异化合作}
\Box{15}{254}{87}{\fontsize{9.8}{14.9}\selectfont\color{ink}广东、台湾和江苏各有三城进入全国前 30；浙江、福建、山东和辽宁各有两城。广东的深圳、广州居第 4、6，浙江的杭州、宁波为第 3、27。省市联读有助于识别核心节点与区域内的不同合作入口。}
\Box{108}{248}{87}{\fontsize{10.1}{14}\selectfont\color{violet}\sffamily \bfseries ④ 地区对标应进一步落到共同任务}
\Box{108}{254}{87}{\fontsize{9.8}{14.9}\selectfont\color{ink}杭州与深圳相差 818.62，广州与成都仅相差 2.98，名次相邻并不意味着分值相近。地区对标需要继续识别供给来源和维护角色：下一节由此进入企业，考察哪些主体能够支撑当地及跨地区的协作。}
\end{canvas}
\begin{canvas}\Header{第三章 · 全球开源影响力排行榜}{violet}\Footer{05}
\Box{15}{27}{180}{\fontsize{20}{27}\selectfont\color{ink}\sffamily \bfseries 3.2 全球开源企业影响力排行榜}
\Box{15}{40}{180}{\fontsize{9.7}{13.5}\selectfont\color{muted}\sffamily 地区中的协作由具体主体承载。企业将仓库供给、开发者参与与长期投入组织起来；全球榜、国别席位及总部例证共同呈现这些主体的布局与影响位置。}
\Box{15}{53}{180}{\includegraphics[width=180mm]{.build/assets/fig-3-04.pdf}}
\Box{15}{199}{180}{\fontsize{7.7}{10.7}\selectfont\color{muted}\sffamily 图 3-4　企业前 20 位与全球 100 强的国别、席位和分值结构。来源：GOSD 2026 企业榜、OpenDigger 复算；前 30 家分值一致。分值窗口为 2025.07—2026.06；国别为组织标签，总部为官网核对例证。}
\Box{15}{214}{87}{\fontsize{10.1}{14}\selectfont\color{violet}\sffamily \bfseries ① 影响、供给规模与参与广度各有重心}
\Box{15}{220}{87}{\fontsize{9.8}{14.9}\selectfont\color{ink}微软、华为、谷歌居前三，OpenRank 分别为 16.71 万、14.13 万、8.25 万。前三家中，华为的仓库最多，微软的关联开发者最多。供给单元、参与基础与网络影响呈现不同组合，需要结合具体项目解释。}
\Box{108}{214}{87}{\fontsize{10.1}{14}\selectfont\color{violet}\sffamily \bfseries ② 100 强覆盖 15 国，总部提供组织入口}
\Box{108}{220}{87}{\fontsize{9.8}{14.9}\selectfont\color{ink}按单一国别标签统计，100 强分布于 15 个国家。地图着色表示企业数量，总部圆点另列六家例证。二者连接组织布局与具体城市；企业项目可以汇聚跨境参与，因此不应把总部国别视为贡献发生地。}
\Box{15}{248}{87}{\fontsize{10.1}{14}\selectfont\color{violet}\sffamily \bfseries ③ 扩展观察范围，显现更完整的企业梯次}
\Box{15}{254}{87}{\fontsize{9.8}{14.9}\selectfont\color{ink}美国占 70 席、中国占 12 席；德国 4 席，荷兰、法国各 2 席，其余十国各 1 席。中国在前 30 名有 4 家，第 31—100 名另有 8 家，涵盖百度、腾讯、PingCAP 等，补充了头部之外的组织观察。}
\Box{108}{248}{87}{\fontsize{10.1}{14}\selectfont\color{violet}\sffamily \bfseries ④ 前五名占 38.5\%，后七十家占 25.3\%}
\Box{108}{254}{87}{\fontsize{9.8}{14.9}\selectfont\color{ink}以 100 强分值合计为分母，前五名占 38.5\%，前十名占 50.2\%。第 31—100 名仍承载 25.3\%，说明观察范围会改变集中度判断。合作资源既可连接大型平台，也应结合具体技术需求识别专业主体。}
\end{canvas}
\begin{canvas}\Header{第三章 · 全球开源影响力排行榜}{violet}\Footer{06}
\Box{15}{27}{180}{\fontsize{20}{27}\selectfont\color{ink}\sffamily \bfseries 3.2.1 企业影响位置的持续与变化}
\Box{15}{40}{180}{\fontsize{9.7}{13.5}\selectfont\color{muted}\sffamily 组织影响既有长期积累，也会随新技术需求发生重排。全球十家企业的年度与月度轨迹，呈现稳定的头部梯次和部分 AI 相关企业的快速前移。}
\Box{15}{53}{180}{\includegraphics[width=180mm]{.build/assets/fig-3-05.pdf}}
\Box{15}{199}{180}{\fontsize{7.7}{10.7}\selectfont\color{muted}\sffamily 图 3-5　全球十家企业的年度与月度位次。来源：GOSD 2026 企业历史榜。年度保留原标签，覆盖窗口待核；月度为 2025.07—2026.06。Anthropic 年度尺度单列，月度按本期前五与后五分组。}
\Box{15}{214}{87}{\fontsize{10.1}{14}\selectfont\color{violet}\sffamily \bfseries ① 长期梯次稳定，部分主体逐步前移}
\Box{15}{220}{87}{\fontsize{9.8}{14.9}\selectfont\color{ink}2022—2026 年度标签中，微软始终第 1，华为由第 5 前移至第 2，英伟达由第 18 前移至第 6。谷歌、Meta、亚马逊保持前五附近。进入前列既可能来自长期积累，也可能表现为连续年度的梯次上移。}
\Box{108}{214}{87}{\fontsize{10.1}{14}\selectfont\color{violet}\sffamily \bfseries ② Anthropic 提供另一条上升路径}
\Box{108}{220}{87}{\fontsize{9.8}{14.9}\selectfont\color{ink}Anthropic 从 2022 年度第 309 位到 2026 年度第 9 位，显著拓宽十家样本的变化幅度。单列尺度可以同时读清其前移与其他头部企业的差别；这一现象提示关注 AI 企业，但仍需项目证据解释原因。}
\Box{15}{248}{87}{\fontsize{10.1}{14}\selectfont\color{violet}\sffamily \bfseries ③ 月度换位改变近期头部组合}
\Box{15}{254}{87}{\fontsize{9.8}{14.9}\selectfont\color{ink}2026 年 3—6 月，华为连续居月榜第 1，微软居第 2。期末英伟达为第 4，Meta 第 5，亚马逊第 9。静态排名之外，连续月份可以显示相对位置何时调整，却不能替代分值或持续维护投入的核查。}
\Box{108}{248}{87}{\fontsize{10.1}{14}\selectfont\color{violet}\sffamily \bfseries ④ 快速前移与阶段波动可以同时发生}
\Box{108}{254}{87}{\fontsize{9.8}{14.9}\selectfont\color{ink}Anthropic 月度位次从期初第 13 前移到 3—4 月第 6，期末回到第 13；Elastic 在 5 月升至第 6，随后为第 8。AI 企业的年度前移值得关注，月度波动则提示继续检验变化的持续性与项目来源。}
\end{canvas}
\begin{canvas}\Header{第三章 · 全球开源影响力排行榜}{violet}\Footer{07}
\Box{15}{27}{180}{\fontsize{20}{27}\selectfont\color{ink}\sffamily \bfseries 3.2.1 企业影响位置的持续与变化：中国}
\Box{15}{40}{180}{\fontsize{9.7}{13.5}\selectfont\color{muted}\sffamily 把同样的时间视角用于中国，可以看见大型平台的稳定位置，也能识别专业技术企业的不同变化。年度轨迹展示梯次形成，月度记录补充近期调整。}
\Box{15}{53}{180}{\includegraphics[width=180mm]{.build/assets/fig-3-06.pdf}}
\Box{15}{199}{180}{\fontsize{7.7}{10.7}\selectfont\color{muted}\sffamily 图 3-6　中国十家企业的年度与月度位次。来源：GOSD 2026 中国企业历史榜。年度覆盖窗口待核；月度为 2025.07—2026.06。按本期前五与后五分组，各面板分别标尺，不代表历年全部企业。}
\Box{15}{214}{87}{\fontsize{10.1}{14}\selectfont\color{violet}\sffamily \bfseries ① 稳定的前两位与逐渐形成的新梯次}
\Box{15}{220}{87}{\fontsize{9.8}{14.9}\selectfont\color{ink}华为、阿里巴巴在五个年度标签中保持第 1、第 2；字节跳动由第 8 前移至第 3，蚂蚁集团与百度在第 3—5 位之间变化。扩大至十家后，长期稳定和内部重排可以同时呈现。}
\Box{108}{214}{87}{\fontsize{10.1}{14}\selectfont\color{violet}\sffamily \bfseries ② 专业主体补充头部之外的观察}
\Box{108}{220}{87}{\fontsize{9.8}{14.9}\selectfont\color{ink}腾讯、PingCAP、飞致云、乐鑫科技与赜睿信息科技在年度样本中形成第 5—12 位的不同轨迹。赜睿信息科技在 2025 年度第 12、2026 年度第 10；这类变化需要结合项目构成，不能仅凭排名解释业务成败。}
\Box{15}{248}{87}{\fontsize{10.1}{14}\selectfont\color{violet}\sffamily \bfseries ③ 近期换位需要与全年位置分开判断}
\Box{15}{254}{87}{\fontsize{9.8}{14.9}\selectfont\color{ink}2026 年 4—6 月，字节跳动连续居中国月榜第 2，阿里巴巴第 3；百度同期为第 6。月度次序补充了年度表，提供进一步查阅项目活动的时间线索，不把最近三个月直接外推为下一年度格局。}
\Box{108}{248}{87}{\fontsize{10.1}{14}\selectfont\color{violet}\sffamily \bfseries ④ 专业主体并不沿同一方向变化}
\Box{108}{254}{87}{\fontsize{9.8}{14.9}\selectfont\color{ink}腾讯在 4—6 月保持第 5；飞致云由 5 月第 10 到 6 月第 8，乐鑫科技同期由第 8 到第 11，赜睿信息科技期末第 9。时序提供跟踪线索，选择合作对象还需核对其技术专长、所在地和维护能力。}
\end{canvas}
\begin{canvas}\Header{第三章 · 全球开源影响力排行榜}{violet}\Footer{08}
\Box{15}{27}{180}{\fontsize{20}{27}\selectfont\color{ink}\sffamily \bfseries 3.2.2 中国企业：空间节点与结构差异}
\Box{15}{40}{180}{\fontsize{9.7}{13.5}\selectfont\color{muted}\sffamily 时间变化之外，中国企业还呈现不同的空间节点和供给结构。将总部、榜内集中度与模型企业的代码侧位置联读，有助于识别地方资源以及跨平台评价的必要性。}
\Box{15}{53}{180}{\includegraphics[width=180mm]{.build/assets/fig-3-07.pdf}}
\Box{15}{199}{180}{\fontsize{7.7}{10.7}\selectfont\color{muted}\sffamily 图 3-7　中国企业排名、总部例证、集中度与模型企业的代码侧位置。来源：GOSD 2026 企业榜；9 家总部据官网于 2026.11.10 核对。集中度限各自前 30 家；人数位次在中国 30 家内重排。}
\Box{15}{214}{87}{\fontsize{10.1}{14}\selectfont\color{violet}\sffamily \bfseries ① 大型组合之外，专业主体构成支撑层}
\Box{15}{220}{87}{\fontsize{9.8}{14.9}\selectfont\color{ink}华为、阿里巴巴、字节跳动、蚂蚁集团居前四，PingCAP、飞致云、乐鑫科技等进入前十。不同企业连接不同规模的仓库和参与者。寻找产业合作伙伴，需要把组织榜单进一步落到具体技术、项目和维护团队。}
\Box{108}{214}{87}{\fontsize{10.1}{14}\selectfont\color{violet}\sffamily \bfseries ② 城市节点提供跨层观察的连接点}
\Box{108}{220}{87}{\fontsize{9.8}{14.9}\selectfont\color{ink}已核对总部例证覆盖北京、杭州、深圳、上海与合肥。杭州连接阿里巴巴与蚂蚁集团，深圳连接华为与腾讯；这些组织节点可与地区榜和项目图谱联读，支持地方识别合作资源，而非直接评价城市整体能力。}
\Box{15}{248}{87}{\fontsize{10.1}{14}\selectfont\color{violet}\sffamily \bfseries ③ 相同前 30 名口径下，榜首权重有别}
\Box{15}{254}{87}{\fontsize{9.8}{14.9}\selectfont\color{ink}为比较相同榜单深度，这里均取各自前 30 家：中国榜前五名占 79.3\%，全球榜占 51.6\%；榜首分别占 54.1\% 和 18.0\%。图 3-4 的 38.5\% 则以全球 100 强为分母，二者不能混用，更不能推导中国的全球份额。}
\Box{108}{248}{87}{\fontsize{10.1}{14}\selectfont\color{violet}\sffamily \bfseries ④ 模型企业需要跨平台评价}
\Box{108}{254}{87}{\fontsize{9.8}{14.9}\selectfont\color{ink}深度求索的代码影响位次为第 26，关联开发者人数位次为第 10；Moonshot AI 分别为第 27、第 16。代码侧位置只揭示一部分协作；本章模型榜与第五章的派生、复用观察将补充其影响路径。}
\end{canvas}
\begin{canvas}\Header{第三章 · 全球开源影响力排行榜}{violet}\Footer{09}
\Box{15}{27}{180}{\fontsize{20}{27}\selectfont\color{ink}\sffamily \bfseries 3.3 全球开源基金会影响力排行榜}
\Box{15}{40}{180}{\fontsize{9.7}{13.5}\selectfont\color{muted}\sffamily 企业投入之外，基金会为跨组织共同建设提供制度与资源支持。将机构分布、项目组合和仓库范围一起观察，才能识别持续维护由谁承载、哪些公共任务值得共同投入。}
\Box{15}{53}{180}{\includegraphics[width=180mm]{.build/assets/fig-3-08.pdf}}
\Box{15}{199}{180}{\fontsize{7.7}{10.7}\selectfont\color{muted}\sffamily 图 3-8　基金会前 20 位、组织国别与项目及仓库覆盖。来源：GOSD 2026 基金会榜、OpenDigger；2026.11.10 标签回溯 2025.07—2026.06。分值沿用可重叠的标签聚合；仓库去重计数；国别补正据官网。}
\Box{15}{214}{87}{\fontsize{10.1}{14}\selectfont\color{violet}\sffamily \bfseries ① 基金会连接不同范围的共同建设网络}
\Box{15}{220}{87}{\fontsize{9.8}{14.9}\selectfont\color{ink}LF、开放原子与 Apache 位居前三。LF 条目包括子基金会，各机构覆盖的技术层和项目组合也不同。排名提供寻找合作网络的入口；支持任务应沿着这些网络进入具体项目，而非按分值直接分配机构绩效。}
\Box{108}{214}{87}{\fontsize{10.1}{14}\selectfont\color{violet}\sffamily \bfseries ② 组织落点集中，协作范围仍可跨越国界}
\Box{108}{220}{87}{\fontsize{9.8}{14.9}\selectfont\color{ink}前二十条目覆盖六个国别，美国占十三，荷兰占三；中国、日本、瑞士与比利时各一。这里呈现机构标签与归属资料，不能据此把其全部项目和贡献者归入同一国家。}
\Box{15}{248}{87}{\fontsize{10.1}{14}\selectfont\color{violet}\sffamily \bfseries ③ 从基金会名称继续进入具体项目}
\Box{15}{254}{87}{\fontsize{9.8}{14.9}\selectfont\color{ink}项目标签揭示组织之内的供给组合。开放原子有四十一个已识别项目，Apache 有二十五个；LF 标签族共识别一百八十一个项目。识别组合有助于把支持资金落实到具体维护对象。}
\Box{108}{248}{87}{\fontsize{10.1}{14}\selectfont\color{violet}\sffamily \bfseries ④ 仓库广度应与项目组织方式一起阅读}
\Box{108}{254}{87}{\fontsize{9.8}{14.9}\selectfont\color{ink}主期可见仓库中，开放原子为 19,043 个，LF 标签族去重为 2,981 个，Apache 为 1,322 个。同一个项目可能拆分为许多仓库，组织标签的覆盖也不均衡，数量差异不能直接解释为产出或公共价值差距。}
\end{canvas}
\begin{canvas}\Header{第三章 · 全球开源影响力排行榜}{violet}\Footer{10}
\Box{15}{27}{180}{\fontsize{20}{27}\selectfont\color{ink}\sffamily \bfseries 3.3.1 LF 内部：多层技术供给与维护重心}
\Box{15}{40}{180}{\fontsize{9.7}{13.5}\selectfont\color{muted}\sffamily Linux 基金会（LF）的子组织覆盖云环境中的软件运行、AI 框架、开发工具与安全等层次。十三个有主期记录的子标签，展示了这些共同建设网络各自的规模与维护重心。}
\Box{15}{53}{180}{\includegraphics[width=180mm]{.build/assets/fig-3-09.pdf}}
\Box{15}{199}{180}{\fontsize{7.7}{10.7}\selectfont\color{muted}\sffamily 图 3-9　LF 子标签的规模、位次与项目组合。来源：OpenDigger；2026.11.10 标签回溯 2025.07—2026.06。跨标签及组织／仓库映射可能重叠，分值不能相加作为 LF 唯一总量；项目数非官方完整托管名录。}
\Box{15}{214}{87}{\fontsize{10.1}{14}\selectfont\color{violet}\sffamily \bfseries ① 云原生与 AI 是两个不同的支撑层}
\Box{15}{220}{87}{\fontsize{9.8}{14.9}\selectfont\color{ink}CNCF 与 PyTorch 分列样本前两位，OpenJS 随后。CNCF 连接云环境中的运行基础，PyTorch 支撑模型计算，OpenJS 连接 Web 开发。智能体与模型应用依赖多层软件协作，支持 AI 也需关注这些已有基础。}
\Box{108}{214}{87}{\fontsize{10.1}{14}\selectfont\color{violet}\sffamily \bfseries ② 安全标签前移，提示持续维护的需求}
\Box{108}{220}{87}{\fontsize{9.8}{14.9}\selectfont\color{ink}OpenSSF 在固定样本中的月度名次由第八升至第六，RISC-V 则由第六至第八，Hyperledger 回到第七。这反映相对协作位置变化，可作为追踪维护重点的信号，不能直接推断技术质量或安全成效。}
\Box{15}{248}{87}{\fontsize{10.1}{14}\selectfont\color{violet}\sffamily \bfseries ③ 核心集中程度提示不同的支持任务}
\Box{15}{254}{87}{\fontsize{9.8}{14.9}\selectfont\color{ink}PyTorch 首位仓库占标签分值的 39.2\%，前三位占 73.5\%；CNCF 分别为 5.2\% 和 10.8\%。前者更应检查核心仓库的维护接续，后者更需要跨项目协调；这些结构差异是核查入口，不等于风险高低的直接结论。}
\Box{108}{248}{87}{\fontsize{10.1}{14}\selectfont\color{violet}\sffamily \bfseries ④ 两个尺度共同确定维护对象}
\Box{108}{254}{87}{\fontsize{9.8}{14.9}\selectfont\color{ink}CNCF 的 70 个项目标签对应 1,339 个主期仓库，PyTorch 的 5 个标签对应 102 个。项目说明建设主题，仓库提供具体工作入口；两者受映射覆盖影响。下一页的 AAIF 则进一步展示新公共组件如何进入基金会。}
\end{canvas}
\begin{canvas}\Header{第三章 · 全球开源影响力排行榜}{violet}\Footer{11}
\Box{15}{27}{180}{\fontsize{20}{27}\selectfont\color{ink}\sffamily \bfseries 3.3.2 AAIF：Agent 协作的公共组件}
\Box{15}{40}{180}{\fontsize{9.7}{13.5}\selectfont\color{muted}\sffamily 智能体是能够协调工具执行任务的 AI 系统。2025 年 12 月成立的 Agentic AI Foundation（AAIF）把相关公共组件纳入共同建设，提供了从组织安排观察协作条件的新案例。}
\Box{15}{53}{180}{\includegraphics[width=180mm]{.build/assets/fig-3-10.pdf}}
\Box{15}{199}{180}{\fontsize{7.7}{10.7}\selectfont\color{muted}\sffamily 图 3-10　AAIF 项目扩展、功能分工、协作记录与治理条件。来源：LF、AAIF 官方公告及 OpenDigger。主期与期后加入分别标注；目录截至 2026.11.10，生命周期政策为 2026.10.15 版。}
\Box{15}{214}{87}{\fontsize{10.1}{14}\selectfont\color{violet}\sffamily \bfseries ① 扩展发生在清楚可辨的时间边界内}
\Box{15}{220}{87}{\fontsize{9.8}{14.9}\selectfont\color{ink}主期内从三个创始项目扩至四个；A2A 与 Agent Router 以及十月的四个新项目均属期后加入。主期末的四个项目与当前十个项目分别陈述，才能看清组织扩展，而不夸大全年托管规模。}
\Box{108}{214}{87}{\fontsize{10.1}{14}\selectfont\color{violet}\sffamily \bfseries ② 中国贡献可以体现为可复用公共组件}
\Box{108}{220}{87}{\fontsize{9.8}{14.9}\selectfont\color{ink}功能地图覆盖指引、互通、执行、流量管理与训练。阿里巴巴贡献的 OpenSandbox 于十月加入，为 AI 应用提供统一沙箱接口。这是中国资源进入全球共同建设的具体例证，实际使用与维护效果仍需继续跟踪。}
\Box{15}{248}{87}{\fontsize{10.1}{14}\selectfont\color{violet}\sffamily \bfseries ③ 协作记录与使用扩散需要分别观察}
\Box{15}{254}{87}{\fontsize{9.8}{14.9}\selectfont\color{ink}十二月至六月，agentgateway 的标签分值由 20.9 至 34.5，MCP 与 goose 的可见分值则下降。协议采用、工具使用和仓库协作并非同一个量；这些变化不能直接解释为市场采用变化或基金会支持成败。}
\Box{108}{248}{87}{\fontsize{10.1}{14}\selectfont\color{violet}\sffamily \bfseries ④ 组织安排要接受实际使用的检验}
\Box{108}{254}{87}{\fontsize{9.8}{14.9}\selectfont\color{ink}现行政策将可运行实现、外部生产使用与持续维护纳入分阶段要求。期后 Gates Foundation 加入并支持非洲开发者参与，体现扩大代表性的组织路径；其公共价值仍需由本地任务、维护接续和服务结果检验。}
\end{canvas}
\begin{canvas}\Header{第三章 · 全球开源影响力排行榜}{violet}\Footer{12}
\Box{15}{27}{180}{\fontsize{20}{27}\selectfont\color{ink}\sffamily \bfseries 3.4 全球开源项目影响力排行榜}
\Box{15}{40}{180}{\fontsize{9.7}{13.5}\selectfont\color{muted}\sffamily 组织怎样转化为具体资源，需要进入项目观察。项目榜将影响位置与技术用途、组织归属、参与规模联读，识别支撑 AI 的基础软件，以及头部之外的专业供给。}
\Box{15}{53}{180}{\includegraphics[width=180mm]{.build/assets/fig-3-11.pdf}}
\Box{15}{199}{180}{\fontsize{7.7}{10.7}\selectfont\color{muted}\sffamily 图 3-11　全球项目前 20 位、百强组织国别与分值结构。来源：GOSD 2026 项目榜、OpenDigger 复算。分值为 2025.07—2026.06；组织字段含发起方与托管方，不代表贡献者来源；参与字段覆盖见附录 A。}
\Box{15}{214}{87}{\fontsize{10.1}{14}\selectfont\color{violet}\sffamily \bfseries ① 头部项目横跨系统、语言与 AI}
\Box{15}{220}{87}{\fontsize{9.8}{14.9}\selectfont\color{ink}OpenHarmony、Azure、C\#/.Net、PyTorch 与 NixOS 位居前五。AI 框架和推理工具已进入头部，但操作系统、语言与云基础设施仍占据重要位置。人工智能的开放发展，需要多层基础软件共同支撑。}
\Box{108}{214}{87}{\fontsize{10.1}{14}\selectfont\color{violet}\sffamily \bfseries ② 组织分布集中，协作边界跨越国别}
\Box{108}{220}{87}{\fontsize{9.8}{14.9}\selectfont\color{ink}百强中，美国标签占 69 项，中国占 7 项，荷兰与德国各 3 项；9 项国别未识别。组织标签适合观察项目承载布局，贡献者却可能来自全球。识别开放协作能力，还须连接项目参与和复用证据。}
\Box{15}{248}{87}{\fontsize{10.1}{14}\selectfont\color{violet}\sffamily \bfseries ③ 中国席位体现系统与 AI 计算支点}
\Box{15}{254}{87}{\fontsize{9.8}{14.9}\selectfont\color{ink}中国 7 个席位中，2 项进入前十、4 项进入前三十，其余位于第 33、37 和 51 位。席位数量与高位节点共同呈现结构位置；据此可以识别合作资源，但不能推导中国对全体全球项目的贡献份额。}
\Box{108}{248}{87}{\fontsize{10.1}{14}\selectfont\color{violet}\sffamily \bfseries ④ 头部与专业支撑层共同构成影响网络}
\Box{108}{254}{87}{\fontsize{9.8}{14.9}\selectfont\color{ink}前十名占百强 OpenRank 合计的 33.6\%，第 31—100 名占 38.7\%。头部项目提供重要支点，较宽的专业层使不同技术任务得到支撑。中国项目页进一步展示这种供给怎样分布在系统、AI 计算和应用工具之间。}
\end{canvas}
\begin{canvas}\Header{第三章 · 全球开源影响力排行榜}{violet}\Footer{13}
\Box{15}{27}{180}{\fontsize{20}{27}\selectfont\color{ink}\sffamily \bfseries 3.4.1 中国项目：从系统支点到 AI 工具}
\Box{15}{40}{180}{\fontsize{9.7}{13.5}\selectfont\color{muted}\sffamily 中国项目呈现系统底座、AI 计算、模型工具与应用框架的分层布局。全球位置显示其已进入哪些协作节点，中国榜进一步揭示企业、基金会与专业项目如何形成供给。}
\Box{15}{53}{180}{\includegraphics[width=180mm]{.build/assets/fig-3-12.pdf}}
\Box{15}{199}{180}{\fontsize{7.7}{10.7}\selectfont\color{muted}\sffamily 图 3-12　中国项目前 20 位、全球百强席位、组织承载与 AI 分层。来源：GOSD 2026 项目榜、OpenDigger 及项目官方文档。全球名次来自同一主期标签榜；技术分层为编辑分类，不跨层加总分值。}
\Box{15}{214}{87}{\fontsize{10.1}{14}\selectfont\color{violet}\sffamily \bfseries ① 头部之外，专业项目已形成梯次}
\Box{15}{220}{87}{\fontsize{9.8}{14.9}\selectfont\color{ink}前七名进入全球百强；Ant Design、openGauss、魔搭、TiDB、verl 等构成其后的专业项目层。榜单不仅用于确认领先者，也可帮助识别设计工具、数据库、模型服务和训练流程中的合作入口。}
\Box{108}{214}{87}{\fontsize{10.1}{14}\selectfont\color{violet}\sffamily \bfseries ② 七个高位节点覆盖三类基础能力}
\Box{108}{220}{87}{\fontsize{9.8}{14.9}\selectfont\color{ink}OpenHarmony 与 openEuler 代表操作系统，CANN、Ascend、MindSpore 与 PaddlePaddle 涉及 AI 计算或框架，Cangjie 对应编程语言。这些资源连接应用运行、模型计算与开发过程，体现多层供给。}
\Box{15}{248}{87}{\fontsize{10.1}{14}\selectfont\color{violet}\sffamily \bfseries ③ 企业与基金会承担不同组织角色}
\Box{15}{254}{87}{\fontsize{9.8}{14.9}\selectfont\color{ink}中国前 30 名中，华为标签对应 6 项，开放原子对应 5 项，字节跳动对应 4 项。企业持续投入与基金会承载共同构成项目发展的组织条件；具体治理开放性仍应考察维护权、决策机制及外部参与。}
\Box{108}{248}{87}{\fontsize{10.1}{14}\selectfont\color{violet}\sffamily \bfseries ④ 中国 AI 供给延伸至训练和应用流程}
\Box{108}{254}{87}{\fontsize{9.8}{14.9}\selectfont\color{ink}按主要用途归类，AI 项目占中国前 30 名的 13 项，涉及算力与框架、模型与后训练、应用与智能体。后训练指在基础模型上继续训练以改善任务表现。分层供给已可识别，实际连接仍需接口与维护证据。}
\end{canvas}
\begin{canvas}\Header{第三章 · 全球开源影响力排行榜}{violet}\Footer{14}
\Box{15}{27}{180}{\fontsize{20}{27}\selectfont\color{ink}\sffamily \bfseries 3.4.2 头部特征：技术、组织与协作规模}
\Box{15}{40}{180}{\fontsize{9.7}{13.5}\selectfont\color{muted}\sffamily 项目名称相近，聚合范围和技术任务仍可能不同。比较技术用途、组织承载、参与规模与对象层次，可以更准确地理解影响位置，并识别下一步需要核查的协作能力。}
\Box{15}{53}{180}{\includegraphics[width=180mm]{.build/assets/fig-3-13.pdf}}
\Box{15}{199}{180}{\fontsize{7.7}{10.7}\selectfont\color{muted}\sffamily 图 3-13　技术用途、组织承载、参与规模与对象层次。来源：GOSD 2026 项目榜及官方文档分类。技术与对象对照各取前 30 名；组织对照为全球百强与中国前 30 名；散点为全球百强。}
\Box{15}{214}{87}{\fontsize{10.1}{14}\selectfont\color{violet}\sffamily \bfseries ① AI 的突出位置须放回全栈结构理解}
\Box{15}{220}{87}{\fontsize{9.8}{14.9}\selectfont\color{ink}全球前 30 名中，AI 与机器学习类有 8 项，中国前 30 名中有 13 项。全球榜同时包含语言、云基础设施、开发工具和系统项目。两榜技术结构不同，提示产业需要将 AI 投入与基础软件维护相互衔接。}
\Box{108}{214}{87}{\fontsize{10.1}{14}\selectfont\color{violet}\sffamily \bfseries ② 公共承载结构是体系能力的一部分}
\Box{108}{220}{87}{\fontsize{9.8}{14.9}\selectfont\color{ink}全球百强中，22 项组织标签指向基金会，4 项指向高校；中国前 30 名中，5 项指向基金会。企业之外的这些载体，为跨组织接续提供制度入口。它们是否形成开放治理，仍须以项目制度与执行记录判断。}
\Box{15}{248}{87}{\fontsize{10.1}{14}\selectfont\color{violet}\sffamily \bfseries ③ 参与规模与协作网络影响并不同步}
\Box{15}{254}{87}{\fontsize{9.8}{14.9}\selectfont\color{ink}Claude Code 的参与账号字段为 32,814，OpenRank 名次为第 25；PyTorch 为 4,806，名次为第 4。二者说明参与规模与网络影响观察的是不同方面。解释差异需回到协作行为与对象范围，不能据此评价效率。}
\Box{108}{248}{87}{\fontsize{10.1}{14}\selectfont\color{violet}\sffamily \bfseries ④ 固定对象边界，才能深入评价协作}
\Box{108}{254}{87}{\fontsize{9.8}{14.9}\selectfont\color{ink}全球前 30 名包括 5 个软件项目、18 个项目生态、7 个项目集合；中国分别为 9、17 和 4。聚合范围不同，比较时必须同时交代仓库清单。下一页通过共同参与和具体接口，把结构观察推进到可验证的连接。}
\end{canvas}
\begin{canvas}\Header{第三章 · 全球开源影响力排行榜}{violet}\Footer{15}
\Box{15}{27}{180}{\fontsize{20}{27}\selectfont\color{ink}\sffamily \bfseries 3.4.3 全球协作：中国项目如何接入与再创造}
\Box{15}{40}{180}{\fontsize{9.7}{13.5}\selectfont\color{muted}\sffamily 排行榜给出位置，关系证据进一步说明连接怎样发生。共同参与账号揭示人员接点，接口与版本适配展示技术接续；中国贡献也由此进入全球再创造的具体过程。}
\Box{15}{53}{180}{\includegraphics[width=180mm]{.build/assets/fig-3-14.pdf}}
\Box{15}{199}{180}{\fontsize{7.7}{10.7}\selectfont\color{muted}\sffamily 图 3-14　跨仓库参与、持续出现与中国项目接口案例。来源：OpenDigger 月度高活跃账号样本、verl 与昇腾 PyTorch 扩展官方文档。账号窗为 2025.07—2026.06；文档版本在主期内；接口支持不等于实际采用。}
\Box{15}{214}{87}{\fontsize{10.1}{14}\selectfont\color{violet}\sffamily \bfseries ① 开放生态存在跨项目的人员接点}
\Box{15}{220}{87}{\fontsize{9.8}{14.9}\selectfont\color{ink}六个仓库的月度高活跃样本中，vLLM 与 SGLang 有 21 个同月重叠账号，PyTorch 与 vLLM 有 16 个。共同账号提供了继续研究协作网络的入口；它本身不证明正式合作、人才迁移或贡献者的国别。}
\Box{108}{214}{87}{\fontsize{10.1}{14}\selectfont\color{violet}\sffamily \bfseries ② 持续协作需要单独观察}
\Box{108}{220}{87}{\fontsize{9.8}{14.9}\selectfont\color{ink}在年度样本中，PyTorch 有 36 个账号在 9—12 个月出现，占 9.8\%；Kubernetes 为 30 个，占 7.0\%。一次性参与与持续出现并存。公共能力建设应跟踪长期维护，但这里的截断样本不能替代全量留存分析。}
\Box{15}{248}{87}{\fontsize{10.1}{14}\selectfont\color{violet}\sffamily \bfseries ③ 中国工具在全球技术基础上再创造}
\Box{15}{254}{87}{\fontsize{9.8}{14.9}\selectfont\color{ink}verl 将 FSDP、Megatron-LM 等训练引擎与 vLLM、SGLang 等生成引擎接入后训练流程。其作用是承接已有资源并组织为可组合工具，让中国项目成为全球技术网络中的连接节点；接口清单尚不等于实际采用规模。}
\Box{108}{248}{87}{\fontsize{10.1}{14}\selectfont\color{violet}\sffamily \bfseries ④ 兼容维护把本土算力接入全球框架}
\Box{108}{254}{87}{\fontsize{9.8}{14.9}\selectfont\color{ink}torch\_npu 连接昇腾与 PyTorch，版本表列出 CANN 9.0.0 的四组配套记录。这类适配与维护使硬件能够使用既有软件框架。下一节将以 CANN 的多仓供给继续观察其组件、参与和版本演进。}
\end{canvas}
\begin{canvas}\Header{第三章 · 全球开源影响力排行榜}{violet}\Footer{16}
\Box{15}{27}{180}{\fontsize{20}{27}\selectfont\color{ink}\sffamily \bfseries 3.5 全球新势力开源仓库与项目排行榜}
\Box{15}{40}{180}{\fontsize{9.7}{13.5}\selectfont\color{muted}\sffamily 从既有项目转向新生供给，首先观察具体仓库。这里的“新”按本期创建记录界定，可能包括迁入或拆分；榜单显示新增协作入口，不能直接作为技术首次出现的时间表。}
\Box{15}{53}{180}{\includegraphics[width=180mm]{.build/assets/fig-3-15.pdf}}
\Box{15}{199}{180}{\fontsize{7.7}{10.7}\selectfont\color{muted}\sffamily 图 3-15　新仓库前 20 位与前 30 位技术及平台结构。来源：GOSD 2026 新仓库榜、逐项用途映射与 OpenDigger。分值窗为 2025.07—2026.06；公开序列复算最大差异 0.02；新建资格的历史映射待核。}
\Box{15}{214}{87}{\fontsize{10.1}{14}\selectfont\color{violet}\sffamily \bfseries ① 头部同时出现应用与计算组件}
\Box{15}{220}{87}{\fontsize{9.8}{14.9}\selectfont\color{ink}OpenClaw 位居首位，随后是 CANN 算子、昇腾框架适配和 MindSpore 仓库。智能体的快速迭代与计算基础软件的开放在同一名单中出现，提示产业应同时关注应用入口和持续支撑它的工程组件。}
\Box{108}{214}{87}{\fontsize{10.1}{14}\selectfont\color{violet}\sffamily \bfseries ② AI 供给的增长有清晰的层次}
\Box{108}{220}{87}{\fontsize{9.8}{14.9}\selectfont\color{ink}30 仓中，AI 基础软件占 18 仓，智能体与工具占 4 仓，合计 73.3\%。算子是执行数学运算的基础组件，框架适配则帮助已有程序使用不同硬件；理解这些分工，才能判断可复用的资源究竟增加在哪里。}
\Box{15}{248}{87}{\fontsize{10.1}{14}\selectfont\color{violet}\sffamily \bfseries ③ 平台让不同协作体系变得可见}
\Box{15}{254}{87}{\fontsize{9.8}{14.9}\selectfont\color{ink}20 个 AtomGit 条目集中于 CANN、Ascend、OpenHarmony 和 MindSpore。多仓库拆分提供了具体参与入口，但 20 个席位并非 20 个独立技术体系；平台分布也不能直接当作国家贡献分布。}
\Box{108}{248}{87}{\fontsize{10.1}{14}\selectfont\color{violet}\sffamily \bfseries ④ 识别新入口，也要识别共同维护对象}
\Box{108}{254}{87}{\fontsize{9.8}{14.9}\selectfont\color{ink}前五仓占 30 仓分值的 43.9\%，其余 25 仓提供较广的建设入口。支持新势力需要把文档、测试与版本协同一起纳入；进一步从仓库聚合为项目，才能辨认哪些组件共同服务于同一套技术体系。}
\end{canvas}
\begin{canvas}\Header{第三章 · 全球开源影响力排行榜}{violet}\Footer{17}
\Box{15}{27}{180}{\fontsize{20}{27}\selectfont\color{ink}\sffamily \bfseries 3.5.1 新势力项目：从组件走向整体供给}
\Box{15}{40}{180}{\fontsize{9.7}{13.5}\selectfont\color{muted}\sffamily 将仓库聚合为项目，可以看见面向使用者的整体供给。与前页组件视角相比，本页突出智能体工作流及多仓体系，展示统计单位与纳入范围如何改变观察结果。}
\Box{15}{53}{180}{\includegraphics[width=180mm]{.build/assets/fig-3-16.pdf}}
\Box{15}{199}{180}{\fontsize{7.7}{10.7}\selectfont\color{muted}\sffamily 图 3-16　新项目前 20 位与前 30 位组织国别、用途和聚合结构。来源：GOSD 2026 新项目榜及逐项映射；分值窗为 2025.07—2026.06。组织国别未识别单列；分值份额以本页 30 项合计为分母。}
\Box{15}{214}{87}{\fontsize{10.1}{14}\selectfont\color{violet}\sffamily \bfseries ① 多仓体系与轻量工具并列出现}
\Box{15}{220}{87}{\fontsize{9.8}{14.9}\selectfont\color{ink}CANN 位列第一，Hermes Agent 与 DeepAgents 随后。CANN 聚合 68 仓，后两项各为单仓，说明新供给既可以来自软件体系的整体开放，也可以来自聚焦任务的工具。分值须与对象范围一同理解。}
\Box{108}{214}{87}{\fontsize{10.1}{14}\selectfont\color{violet}\sffamily \bfseries ② 中国贡献覆盖底座与应用工具}
\Box{108}{220}{87}{\fontsize{9.8}{14.9}\selectfont\color{ink}已识别的中国组织有 7 项，包含 CANN、OpenViking、Astron、Nanobot、QwenPaw 等，跨越计算支撑、企业工具与高校研发。另有 17 项组织国别未识别，因此当前地图呈现的是已知布局，不能推断完整全球格局。}
\Box{15}{248}{87}{\fontsize{10.1}{14}\selectfont\color{violet}\sffamily \bfseries ③ 项目入口更集中于智能体工作流}
\Box{15}{254}{87}{\fontsize{9.8}{14.9}\selectfont\color{ink}25 项主要用于智能体与相关工具，占 83.3\%；AI 基础软件为 4 项。与仓库榜的组件结构对照，可见从计算支撑到任务执行的两层供给。两榜平台与纳入范围不同，这一差异也不能全部归因于仓库聚合。}
\Box{108}{248}{87}{\fontsize{10.1}{14}\selectfont\color{violet}\sffamily \bfseries ④ 多仓聚合让体系供给变得可见}
\Box{108}{254}{87}{\fontsize{9.8}{14.9}\selectfont\color{ink}多仓项目为 10 项，占项目数三分之一，却占榜内分值的 67.5\%。这提示评价新势力需兼看组件范围、接口配套和跨仓协作。仓库多并不自动说明体系成熟，CANN 为核对这些连接提供了具体案例。}
\end{canvas}
\begin{canvas}\Header{第三章 · 全球开源影响力排行榜}{violet}\Footer{18}
\Box{15}{27}{180}{\fontsize{20}{27}\selectfont\color{ink}\sffamily \bfseries 3.5.2 CANN：把多仓供给组织为体系能力}
\Box{15}{40}{180}{\fontsize{9.7}{13.5}\selectfont\color{muted}\sffamily CANN（神经网络计算架构）是华为连接 AI 框架与昇腾处理器的异构计算软件栈。其多仓结构提供了具体案例：开放组件如何通过接口、版本协同与社区维护形成体系能力。}
\Box{15}{53}{180}{\includegraphics[width=180mm]{.build/assets/fig-3-17.pdf}}
\Box{15}{199}{180}{\fontsize{7.7}{10.7}\selectfont\color{muted}\sffamily 图 3-17　CANN 的 68 仓画像及期后延伸。来源：OpenDigger、CANN 官方首页与华为开放及治理公告。主期为 2025.07—2026.06；曲线期后段为 2026.07—10；主页设施为 11-10 观察。仓库按 ID 去重。}
\Box{15}{214}{87}{\fontsize{10.1}{14}\selectfont\color{violet}\sffamily \bfseries ① 完整范围使结构性贡献可被识别}
\Box{15}{220}{87}{\fontsize{9.8}{14.9}\selectfont\color{ink}68 仓汇总分值与原榜一致，前 12 仓占 77.4\%。Transformer 与神经网络算子居前，编程、通信及推理实践也进入头部。中国 AI 开源的贡献因而可以落实到可协作的软件层次，而不止停留在模型名称。}
\Box{108}{214}{87}{\fontsize{10.1}{14}\selectfont\color{violet}\sffamily \bfseries ② 分层开放为上游协作提供接口}
\Box{108}{220}{87}{\fontsize{9.8}{14.9}\selectfont\color{ink}官方资料将加速库、编程编译、图引擎和运行时等分层呈现，并提出与 PyTorch、vLLM 等社区协作。接口开放为适配与性能改进提供条件；是否形成跨境复用，还应继续核对版本兼容、上游接纳和实际采用。}
\Box{15}{248}{87}{\fontsize{10.1}{14}\selectfont\color{violet}\sffamily \bfseries ③ 影响力增长伴随组件范围扩展}
\Box{15}{254}{87}{\fontsize{9.8}{14.9}\selectfont\color{ink}主期末 6 月月度 OpenRank 为 3,520，10 月升至 4,856；有影响力记录的仓库由 63 个增至 90 个。增长体现协作网络与组件范围的共同扩展，不能直接解释为单仓效率、代码质量或商业采用提高。}
\Box{108}{248}{87}{\fontsize{10.1}{14}\selectfont\color{violet}\sffamily \bfseries ④ 共同参与要逐步沉淀为长期责任}
\Box{108}{254}{87}{\fontsize{9.8}{14.9}\selectfont\color{ink}主期 5,282 个去机器人协作账号中，1,868 个跨仓参与，1,869 个至少在三个月出现。免费算力、示例和评测提供参与条件；后续应追踪兼容维护、适配成本与真实复用，检验这些投入能否降低使用门槛。}
\end{canvas}
\begin{canvas}\Header{第三章 · 全球开源影响力排行榜}{violet}\Footer{19}
\Box{15}{27}{180}{\fontsize{20}{27}\selectfont\color{ink}\sffamily \bfseries 3.6 全球开源模型影响力排行榜}
\Box{15}{40}{180}{\fontsize{9.7}{13.5}\selectfont\color{muted}\sffamily 代码项目之外，模型通过派生、适配和复用形成另一类影响网络。ModelRank 用这些关系观察网络位置；2026 年 6 月百强快照展示的是复用节点，而非模型能力或开放程度认证。}
\Box{15}{53}{180}{\includegraphics[width=180mm]{.build/assets/fig-3-18.pdf}}
\Box{15}{199}{180}{\fontsize{7.7}{10.7}\selectfont\color{muted}\sffamily 图 3-18　全球模型前 20 位、发布组织分布与百强结构。来源：ModelRank 2026 年 6 月归档，非全年累计；组织据官方资料于 11 月核验。94 席可定位，6 席未知；地图不表示贡献者、训练数据或用户所在地。}
\Box{15}{214}{87}{\fontsize{10.1}{14}\selectfont\color{violet}\sffamily \bfseries ① 头部影响位置跨越语言、视觉与机器人}
\Box{15}{220}{87}{\fontsize{9.8}{14.9}\selectfont\color{ink}FLUX.1-dev 和 Stable Diffusion XL 基础版居前两位；前十有 5 个文本生成版本、4 个图像生成版本和 1 个机器人模型。榜单呈现的是多种复用节点，公共资源建设应兼顾这些任务差异。}
\Box{108}{214}{87}{\fontsize{10.1}{14}\selectfont\color{violet}\sffamily \bfseries ② 地图呈现组织布局，贡献仍跨越国界}
\Box{108}{220}{87}{\fontsize{9.8}{14.9}\selectfont\color{ink}按发布组织归属，美国占 50 席、中国占 37 席，其余 7 席分布于 6 国，另有 6 席未知。组织所在国提供空间入口；底座继承、社区适配和跨国维护关系，仍须沿具体模型继续追踪。}
\Box{15}{248}{87}{\fontsize{10.1}{14}\selectfont\color{violet}\sffamily \bfseries ③ 席位数量与头部位置并不同序}
\Box{15}{254}{87}{\fontsize{9.8}{14.9}\selectfont\color{ink}中国的 37 席中有 16 席进入前三十；美国的 50 席中有 10 席进入前三十。英德虽各仅一席，却占据前两位。数量、排名梯次与技术专长共同构成格局，任何单项都不足以概括国家能力。}
\Box{108}{248}{87}{\fontsize{10.1}{14}\selectfont\color{violet}\sffamily \bfseries ④ 全景观察让支撑性任务显现}
\Box{108}{254}{87}{\fontsize{9.8}{14.9}\selectfont\color{ink}文本生成占 38 席，其余版本还覆盖图文理解、语音、检索与特征表示，5 席无任务标注。面向产业与公共服务，供给广度应落实到任务适用性；中国模型页据此进一步观察通用系列和科研专长。}
\end{canvas}
\begin{canvas}\Header{第三章 · 全球开源影响力排行榜}{violet}\Footer{20}
\Box{15}{27}{180}{\fontsize{20}{27}\selectfont\color{ink}\sffamily \bfseries 3.6.1 中国模型：榜内位置与供给结构}
\Box{15}{40}{180}{\fontsize{9.7}{13.5}\selectfont\color{muted}\sffamily 全球百强中可识别的 37 个中国发布模型，连接了通用系列与科研专长。本页列出其中前 20 位，观察其任务、型号规模和组织节点；这是百强子集，不是完整中国模型总榜。}
\Box{15}{53}{180}{\includegraphics[width=180mm]{.build/assets/fig-3-19.pdf}}
\Box{15}{199}{180}{\fontsize{7.7}{10.7}\selectfont\color{muted}\sffamily 图 3-19　中国发布模型的全球位次、组织节点、任务与规模。来源：ModelRank 6 月归档、组织资料及 bge-m3 模型卡。B 表示十亿参数；型号规模不等于运行成本。Qwen 36 席中，31 席标称规模不超过 10B。}
\Box{15}{214}{87}{\fontsize{10.1}{14}\selectfont\color{violet}\sffamily \bfseries ① 多代模型共同构成持续可复用的底座}
\Box{15}{220}{87}{\fontsize{9.8}{14.9}\selectfont\color{ink}Qwen 占 36 席，9 席进入全球前二十；榜内前二十的中国版本均来自 Qwen。Qwen2.5、Qwen3、Qwen3.5 与 Qwen3.6 同时出现，说明既有版本仍占据复用位置，不能只追踪最新发布。}
\Box{108}{214}{87}{\fontsize{10.1}{14}\selectfont\color{violet}\sffamily \bfseries ② 企业系列与科研专长形成互补}
\Box{108}{220}{87}{\fontsize{9.8}{14.9}\selectfont\color{ink}样本连接杭州的阿里巴巴与北京智源两类主体：前者提供通用模型系列，后者以 bge-m3 进入第 51 位。组织节点揭示供给的不同载体，尚不足以推出城市整体实力、人才聚集或政策成效。}
\Box{15}{248}{87}{\fontsize{10.1}{14}\selectfont\color{violet}\sffamily \bfseries ③ 生成之外，检索能力也进入头部网络}
\Box{15}{254}{87}{\fontsize{9.8}{14.9}\selectfont\color{ink}37 席包括文本生成 22 席、图文输入转文本 13 席，以及特征提取、句子相似度各 1 席。检索与表示模型有助于把知识资源接入应用；单看对话模型会遗漏这类支撑环节。}
\Box{108}{248}{87}{\fontsize{10.1}{14}\selectfont\color{violet}\sffamily \bfseries ④ 多种规模为不同部署条件提供选择}
\Box{108}{254}{87}{\fontsize{9.8}{14.9}\selectfont\color{ink}Qwen 的 36 席中，31 席的标称参数量不超过 10B，含 7 席不超过 1B。这是资源可及性的供给线索；实际成本还受结构、精度与硬件影响。与 DeepSeek 家族联读，可以进一步看见不同使用条件下的互补。}
\end{canvas}
\begin{canvas}\Header{第三章 · 全球开源影响力排行榜}{violet}\Footer{21}
\Box{15}{27}{180}{\fontsize{20}{27}\selectfont\color{ink}\sffamily \bfseries 3.6.2 Qwen 与 DeepSeek：互补与接续}
\Box{15}{40}{180}{\fontsize{9.7}{13.5}\selectfont\color{muted}\sffamily 从单个版本转向家族与接续者，Qwen 的社区复用和 DeepSeek 的大型模型使用呈现不同路径。结合模型卡中的来源关系，可以进一步识别基础权重、推理经验与适配维护各自的贡献。}
\Box{15}{53}{180}{\includegraphics[width=180mm]{.build/assets/fig-3-20.pdf}}
\Box{15}{199}{180}{\fontsize{7.7}{10.7}\selectfont\color{muted}\sffamily 图 3-20　Qwen 与 DeepSeek 的采用结构、社区转换与跨家族接续。来源：ATOM（2026.04）、HF 夏季报告（8 月固定版）、DeepSeek 与 Unsloth 模型卡及 6 月榜单。窗口、平台和分母各自标明，不作统一份额。}
\Box{15}{214}{87}{\fontsize{10.1}{14}\selectfont\color{violet}\sffamily \bfseries ① 两种优势对应不同的使用条件}
\Box{15}{220}{87}{\fontsize{9.8}{14.9}\selectfont\color{ink}ATOM 样本中，Qwen 占小于 10B 组累计下载的 44\%，DeepSeek 占 250B 以上组的 47\%；调用数据补充了后者的服务使用。部署与复用路径不同，不能由六月榜内未见 DeepSeek 就推断其缺乏影响。}
\Box{108}{214}{87}{\fontsize{10.1}{14}\selectfont\color{violet}\sffamily \bfseries ② 社区承担把模型变得可用的工作}
\Box{108}{220}{87}{\fontsize{9.8}{14.9}\selectfont\color{ink}GGUF 是便于本地推理的模型文件格式。HF 的 Qwen GGUF 仓库观察中，非 Qwen 账户承担 99.8\% 的发布，显示适配分发的社区作用。转换数量尚不证明质量或部署，维护与复核仍是资源可用的必要条件。}
\Box{15}{248}{87}{\fontsize{10.1}{14}\selectfont\color{violet}\sffamily \bfseries ③ 中国贡献体现为底座与推理经验}
\Box{15}{254}{87}{\fontsize{9.8}{14.9}\selectfont\color{ink}通过蒸馏，DeepSeek-R1 将生成的推理示例用于训练较小模型，形成 Qwen 与 Llama 底座上的六个版本，后续又接续到 Qwen3。基础权重与训练样本分别贡献能力，体现跨家族、跨主体再创造。}
\Box{108}{248}{87}{\fontsize{10.1}{14}\selectfont\color{violet}\sffamily \bfseries ④ 影响评价应同时看供给者与接续者}
\Box{108}{254}{87}{\fontsize{9.8}{14.9}\selectfont\color{ink}Qwen 与 Unsloth 的同名版本分别进入第 6、第 10 位，说明基础供给与适配分发均可能获得网络位置。从本章的主体识别进入第四、五章，需要继续追踪接口、派生来源、开放条件与共同维护责任。}
\end{canvas}
\begin{canvas}\Header{第三章 · 全球开源影响力排行榜}{violet}\Footer{22}
\Box{15}{27}{180}{\fontsize{20}{27}\selectfont\color{ink}\sffamily \bfseries 本章小结　从影响位置到共同建设}
\Box{15}{44}{180}{\fontsize{10.3}{16}\selectfont\color{ink}六类对象的联读，将影响力从单一名次扩展为地域、组织、资源与接续关系。由此形成的行动方向，是把可见的网络位置转化为具体共同任务，并以维护与使用结果持续检验。}
\Box{15}{67}{180}{\fontsize{11.3}{16}\selectfont\color{violet}\sffamily \bfseries 三项判断}
\Box{15}{79}{180}{\fontsize{10.7}{15.5}\selectfont\color{ink}\sffamily \bfseries 地域规模、组织供给与网络影响需要同时观察}
\Box{15}{85.5}{180}{\fontsize{9.7}{14.5}\selectfont\color{ink}地区的两类位次并不同序；企业与项目的聚合范围各异。分层识别对象，才能找到适合不同地方与产业的合作入口。}
\Box{15}{102}{180}{\fontsize{10.7}{15.5}\selectfont\color{ink}\sffamily \bfseries 中国贡献体现为多层资源与全球接续}
\Box{15}{108.5}{180}{\fontsize{9.7}{14.5}\selectfont\color{ink}CANN、verl、Qwen 与 DeepSeek 分别展示计算支撑、训练工具、模型底座与推理经验。价值体现在具体资源及其连接，不能用单一榜单概括全面领先。}
\Box{15}{125}{180}{\fontsize{10.7}{15.5}\selectfont\color{ink}\sffamily \bfseries 公共能力需要有组织、可持续的共同维护}
\Box{15}{131.5}{180}{\fontsize{9.7}{14.5}\selectfont\color{ink}基金会项目组合、跨仓参与与社区适配提供了组织和工作入口。长期价值仍取决于责任、运行条件及实际使用，不能由数量与热度自动推出。}
\Box{15}{151}{180}{\fontsize{11.3}{16}\selectfont\color{violet}\sffamily \bfseries 行动启示｜未来 6—12 个月的可检验任务}
\Box{15}{163}{47}{\fontsize{10}{14}\selectfont\color{violet}\sffamily \bfseries 把地区优势落实到共同项目}
\Box{15}{173}{47}{\fontsize{8.2}{11.5}\selectfont\color{muted}\sffamily 地方部门、高校与企业}
\Box{66}{163}{129}{\fontsize{9.4}{14.2}\selectfont\color{ink}据图 3-1—3-7，联读地域、组织与项目资源。配置协调人员和维护岗位，选择有需求的跨主体任务；以被接收的贡献、共同交付及持续责任检验进展。}
\Rule{15}{186}{195}{line}
\Box{15}{188}{47}{\fontsize{10}{14}\selectfont\color{violet}\sffamily \bfseries 共同投入兼容与维护基础}
\Box{15}{198}{47}{\fontsize{8.2}{11.5}\selectfont\color{muted}\sffamily 企业、基金会与公共技术平台}
\Box{66}{188}{129}{\fontsize{9.4}{14.2}\selectfont\color{ink}据图 3-9、3-14、3-17，围绕明确接口投入设备、测试与维护人员。观察兼容覆盖、问题响应和版本接续，避免以仓库数量替代建设成效。}
\Rule{15}{211}{195}{line}
\Box{15}{213}{47}{\fontsize{10}{14}\selectfont\color{violet}\sffamily \bfseries 让模型接续可追溯、可复用}
\Box{15}{223}{47}{\fontsize{8.2}{11.5}\selectfont\color{muted}\sffamily 发布团队、平台与接续社区}
\Box{66}{213}{129}{\fontsize{9.4}{14.2}\selectfont\color{ink}据图 3-18—3-20，保留来源、版本、转换说明与许可条件；支持适配和独立验证。记录实际采用、上游接纳与维护交接，兼顾基础供给者和接续者。}
\Rule{15}{236}{195}{line}
\Box{15}{238}{47}{\fontsize{10}{14}\selectfont\color{violet}\sffamily \bfseries 由本地伙伴定义普惠任务}
\Box{15}{248}{47}{\fontsize{8.2}{11.5}\selectfont\color{muted}\sffamily 本地公共机构、社区与国际伙伴}
\Box{66}{238}{129}{\fontsize{9.4}{14.2}\selectfont\color{ink}据图 3-1、3-10、3-19，本地团队先确定语言与场景需求，伙伴提供运行和维护预算。以成本、本地效果、独立运维及治理参与检验，不把加入组织当作成效。}
\Rule{15}{261}{195}{line}
\Box{15}{266}{180}{\fontsize{9}{13.1}\selectfont\color{muted}上述为证据引出的行动建议，尚非已验证效果。第四、五章继续检验技术连接、模型复用与开放条件；第六章综合跨章证据，形成完整的战略建议与实施安排。}
\end{canvas}
\begin{canvas}\Header{附录 A · 第三章方法与证据}{violet}\Footer{A1}
\Box{15}{27}{180}{\fontsize{20}{27}\selectfont\color{ink}\sffamily \bfseries 附录 A　指标与时间口径}
\Box{15}{43}{180}{\fontsize{9.9}{16}\selectfont\color{ink}本章同时观察地区、企业、基金会、项目、仓库和模型。对象、聚合范围与时间决定指标的含义；以下说明适用于全章，特殊限制仍在对应图题与正文保留。}
\Box{15}{64}{180}{\fontsize{11.3}{16}\selectfont\color{violet}\sffamily \bfseries A1.1　核心指标}
\Box{15}{77}{28}{\fontsize{10.2}{15}\selectfont\color{violet}\sffamily \bfseries 网络影响}
\Box{47}{77}{148}{\fontsize{9.6}{14.8}\selectfont\color{ink}OpenRank 根据协作关系及相连节点的重要性计算影响。企业与项目主榜按主期月度分值汇总；不同对象及重叠标签的分值不相加。模型使用独立的 ModelRank 月榜。}
\Box{15}{102}{28}{\fontsize{10.2}{15}\selectfont\color{violet}\sffamily \bfseries 参与与供给}
\Box{47}{102}{148}{\fontsize{9.6}{14.8}\selectfont\color{ink}关联开发者、参与账号和仓库数沿用各表字段，分别描述参与基础和代码单元；不等于企业员工、职业人口、原创作品或实际用户。部分参与字段的去重和覆盖尚待核。}
\Box{15}{127}{28}{\fontsize{10.2}{15}\selectfont\color{violet}\sffamily \bfseries 席位与份额}
\Box{47}{127}{148}{\fontsize{9.6}{14.8}\selectfont\color{ink}席位是榜单中符合条件的对象数。分值份额以注明的前 30 名或前 100 名合计为分母；中国子榜、全球榜与不同深度的比例不能互换。}
\Box{15}{152}{28}{\fontsize{10.2}{15}\selectfont\color{violet}\sffamily \bfseries 地域与组织}
\Box{47}{152}{148}{\fontsize{9.6}{14.8}\selectfont\color{ink}地区榜按原地域字段；企业按组织国别和总部例证，项目按发起／托管标签，模型按发布组织。它们都不代表全部贡献者、训练数据或用户所在地。未知保持单列。}
\Box{15}{177}{28}{\fontsize{10.2}{15}\selectfont\color{violet}\sffamily \bfseries 复用与公共价值}
\Box{47}{177}{148}{\fontsize{9.6}{14.8}\selectfont\color{ink}模型派生、格式转换、下载和调用记录反映不同接续与使用路径，不等于同一个用户量。接口支持和组织安排提供条件；公共成效仍需成本、本地化、维护和结果证据。}
\Box{15}{207}{180}{\fontsize{11.3}{16}\selectfont\color{violet}\sffamily \bfseries A1.2　三类时间分开阅读}
\Box{15}{219}{180}{\fontsize{9.6}{14.8}\selectfont\color{ink}主期及期末：报告主期为 2025-07-01—2026-06-30；ModelRank 为期末 6 月单月快照。地区原榜统计起止及企业历史年度标签的具体覆盖仍待核，不默认其均与主期一致。}
\Box{15}{241}{180}{\fontsize{9.6}{14.8}\selectfont\color{ink}历史与外部窗口：企业 2022—2026 年度轨迹保留原标签；ATOM 使用其研究样本和观察窗。HF 夏季报告分析 1—7 月活动并于 8 月补充发布，不能回填至 6 月榜单。}
\Box{15}{263}{180}{\fontsize{9.3}{14}\selectfont\color{muted}期后观察：CANN 7—10 月走势、AAIF 后续加入与 10 月政策、11 月组织和设施核验各自标注。当前标签回溯主期，不证明历史映射始终相同。}
\end{canvas}
\begin{canvas}\Header{附录 A · 第三章方法与证据}{violet}\Footer{A2}
\Box{15}{27}{180}{\fontsize{20}{27}\selectfont\color{ink}\sffamily \bfseries 附录 A　对象边界与核验范围}
\Box{15}{45}{180}{\fontsize{11.3}{16}\selectfont\color{violet}\sffamily \bfseries A2.1　地区与企业}
\Box{15}{53}{180}{\fontsize{9.4}{14.4}\selectfont\color{ink}行政区榜混合州、省、构成国等单元；地区参与人数的单位与覆盖待核，因此采用同一 30 地区的位次比较。企业全球前 30 家分值与 OpenDigger 主期重算一致；总部图为九家中国企业与六家全球企业的例证，不能代表所有活动发生地。}
\Box{15}{80}{180}{\fontsize{11.3}{16}\selectfont\color{violet}\sffamily \bfseries A2.2　基金会标签与去重}
\Box{15}{88}{180}{\fontsize{9.4}{14.4}\selectfont\color{ink}LF 父子标签及组织／仓库映射可能重叠。主榜保留原标签聚合：LF 196,128.27、开放原子 99,965.35、Apache 81,247.51；按平台、仓库、月份去重后分别为 165,911.86、99,920.91、59,322.07。两套值不可择取混用；正文仓库计数已去重，项目数为识别到的标签数。}
\Box{15}{121}{180}{\fontsize{11.3}{16}\selectfont\color{violet}\sffamily \bfseries A2.3　项目分类与参与样本}
\Box{15}{129}{180}{\fontsize{9.4}{14.4}\selectfont\color{ink}技术用途按官方文档作单一主要用途分类，不覆盖全部 AI 相关性。项目对象分为软件、生态与集合；全球百强和中国前三十分值复算一致，但参与字段的机器人与跨平台身份处理仍待核。六仓库关系图每仓每月至多取 100 个高活跃账号，剔除已识别机器人后去重；同月重叠不是正式合作，持续出现不是全量留存。}
\Box{15}{167}{180}{\fontsize{11.3}{16}\selectfont\color{violet}\sffamily \bfseries A2.4　新势力与 CANN}
\Box{15}{175}{180}{\fontsize{9.4}{14.4}\selectfont\color{ink}“新”来自创建记录，迁入、拆分或改名可能影响资格。CANN 按仓库 ID 合并三组改名后，主期 68 仓共 14,719.47 分；去机器人协作账号 5,282 个，区别于榜表参与字段 5,460 个。6 月 63 仓、10 月 90 仓为当月有影响分值记录的仓库，不能与主期并集直接相比。}
\Box{15}{208}{180}{\fontsize{11.3}{16}\selectfont\color{violet}\sffamily \bfseries A2.5　模型与开放条件}
\Box{15}{216}{180}{\fontsize{9.4}{14.4}\selectfont\color{ink}6 月百强保留 22 个发布账户；同一机构关联账户在地域汇总时合并。6 席国别、5 席任务未知；ModelRank 的完整算法与历史版本链待进一步复核。Qwen 规模按型号标称总参数量分组，35B-A3B 计 35B。DeepSeek 不在本期百强，外部采用证据另列，不补排。}
\Box{15}{254}{180}{\fontsize{9.4}{14.3}\selectfont\color{muted}入榜不等于完整开源认证。公开仓库不必然开放全部产品，模型权重可得也不自动满足全部开源条件；逐项许可与实际部署效果仍需核查。来源记录、数据及方法保留在结构化源稿中。}
\end{canvas}
\begin{canvas}\Header{附录 A · 第三章方法与证据}{violet}\Footer{A3}
\Box{15}{27}{180}{\fontsize{20}{27}\selectfont\color{ink}\sffamily \bfseries 附录 A　数据与技术证据}
\Box{15}{43}{180}{\fontsize{9.5}{14.5}\selectfont\color{muted}\sffamily 以下保留必要的原始入口；正文中的简明来源与这些资料一一对应。报告榜单取固定归档，OpenDigger 数据复核及组织核验日期为 2026-11-10。}
\Box{15}{63.0}{180}{\fontsize{9.5}{14}\selectfont\color{ink}\sffamily \bfseries \href{https://github.com/X-lab2017/2026-Global-Open-Source-Development-Report/blob/c51e3441361407b68f10b339d97eff288319b503/data/processed/chapter3}{GOSD 2026 行政区划、企业与基金会榜}}
\Box{15}{69.0}{180}{\fontsize{8.8}{12.8}\selectfont\color{muted}\sffamily 图 3-1—3-10；国别补正据 \href{https://duckdb.foundation/}{DuckDB 登记资料}与 \href{https://www.eclipse.org/org/foundation/annual-reports/2025/}{Eclipse 2025 年报}。}
\Box{15}{78.8}{180}{\fontsize{9.5}{14}\selectfont\color{ink}\sffamily \bfseries \href{https://github.com/X-lab2017/2026-Global-Open-Source-Development-Report/blob/c51e3441361407b68f10b339d97eff288319b503/figures/chapter3/leaderboards/projects.html}{GOSD 2026 全球项目百强及中国项目榜}}
\Box{15}{84.8}{180}{\fontsize{8.8}{12.8}\selectfont\color{muted}\sffamily 图 3-11—3-14；全球百强来自项目榜归档，中国前 30 位为同版数据。}
\Box{15}{94.6}{180}{\fontsize{9.5}{14}\selectfont\color{ink}\sffamily \bfseries \href{https://github.com/X-lab2017/2026-Global-Open-Source-Development-Report/blob/c51e3441361407b68f10b339d97eff288319b503/preview/content/gosr/emerging-reference.json}{GOSD 2026 新仓库、新项目榜及用途映射}}
\Box{15}{100.6}{180}{\fontsize{8.8}{12.8}\selectfont\color{muted}\sffamily 图 3-15—3-17；两榜各 30 个对象，逐项保留组织、用途与官方入口。}
\Box{15}{110.4}{180}{\fontsize{9.5}{14}\selectfont\color{ink}\sffamily \bfseries \href{https://github.com/X-lab2017/2026-Global-Open-Source-Development-Report/blob/c51e3441361407b68f10b339d97eff288319b503/preview/content/gosr/model-monthly-source/modelrank-2026-06-top100.raw.json}{ModelRank 2026 年 6 月 Top 100 原始归档}}
\Box{15}{116.4}{180}{\fontsize{8.8}{12.8}\selectfont\color{muted}\sffamily 图 3-18—3-20；逐项保留名称、分值、任务与发布账户。}
\Box{15}{126.2}{180}{\fontsize{9.5}{14}\selectfont\color{ink}\sffamily \bfseries \href{https://github.com/X-lab2017/open-digger/blob/5c2ffc4f5b3dd37ebcb20514c7abfee6853bd830/src/metrics/indices.ts}{OpenDigger：指标实现与标签体系}}
\Box{15}{132.2}{180}{\fontsize{8.8}{12.8}\selectfont\color{muted}\sffamily 企业与项目复算、基金会标签回溯及 CANN 查询；标签与 SQL 明细随源稿保存。}
\Box{15}{142.0}{180}{\fontsize{9.5}{14}\selectfont\color{ink}\sffamily \bfseries \href{https://oss.open-digger.cn/github/pytorch/pytorch/activity_details.json}{OpenDigger：月度高活跃账号样本}}
\Box{15}{148.0}{180}{\fontsize{8.8}{12.8}\selectfont\color{muted}\sffamily 图 3-14；PyTorch、vLLM、SGLang、LangChain、Dify、Kubernetes，共六仓。}
\Box{15}{157.8}{180}{\fontsize{9.5}{14}\selectfont\color{ink}\sffamily \bfseries \href{https://github.com/verl-project/verl/blob/38b57931f0c4f4da1ff85a7d54c2d8bbc239163d/README.md}{verl 官方说明，2026-05-29 固定版本}}
\Box{15}{163.8}{180}{\fontsize{8.8}{12.8}\selectfont\color{muted}\sffamily 图 3-14：训练、生成引擎与模型、硬件支持。}
\Box{15}{173.60000000000002}{180}{\fontsize{9.5}{14}\selectfont\color{ink}\sffamily \bfseries \href{https://github.com/Ascend/pytorch/blob/554f7cd18e52dc2f94a5715c100b7d15b7b65de0/README.md}{昇腾 PyTorch 扩展说明，2026-06-16 固定版本}}
\Box{15}{179.60000000000002}{180}{\fontsize{8.8}{12.8}\selectfont\color{muted}\sffamily 图 3-14：扩展定位、CANN 与 PyTorch 的版本配套。}
\Box{15}{194}{180}{\fontsize{11.3}{16}\selectfont\color{violet}\sffamily \bfseries 发布组织与地域核验}
\Box{15}{205}{180}{\fontsize{9.2}{13.8}\selectfont\color{ink}地图表示组织所在地或榜表标签。模型百强可识别 94 席，以下按组织合并关联账户；未知为 intfloat 4 席、KBlueLeaf 1 席、cmz1024 1 席，不按姓名或语言推断国别。}
\Box{15}{225.0}{87}{\fontsize{8.7}{12}\selectfont\color{muted}\sffamily \href{https://about.google/company-info/locations/}{Google · 美国 20}}
\Box{15}{231.6}{87}{\fontsize{8.7}{12}\selectfont\color{muted}\sffamily \href{https://www.sec.gov/Archives/edgar/data/1326801/000162828026003942/meta-20251231.htm}{Meta · 美国 13}}
\Box{15}{238.2}{87}{\fontsize{8.7}{12}\selectfont\color{muted}\sffamily \href{https://openai.com/policies/terms-of-use/}{OpenAI · 美国 7}}
\Box{15}{244.8}{87}{\fontsize{8.7}{12}\selectfont\color{muted}\sffamily \href{https://huggingface.co/terms-of-service}{Hugging Face · 美国 2}}
\Box{15}{251.4}{87}{\fontsize{8.7}{12}\selectfont\color{muted}\sffamily \href{https://www.clsp.jhu.edu/contact/}{JHU CLSP · 美国 2}}
\Box{15}{258.0}{87}{\fontsize{8.7}{12}\selectfont\color{muted}\sffamily \href{https://www.krea.ai/blog/new-krea}{Krea · 美国 2}}
\Box{15}{264.6}{87}{\fontsize{8.7}{12}\selectfont\color{muted}\sffamily \href{https://www.liquid.ai/privacy-policy}{Liquid AI · 美国 2}}
\Box{15}{271.2}{87}{\fontsize{8.7}{12}\selectfont\color{muted}\sffamily \href{https://www.ycombinator.com/companies/unsloth-ai}{Unsloth · 美国 2}}
\Box{108}{225.0}{87}{\fontsize{8.7}{12}\selectfont\color{muted}\sffamily \href{https://home.alibabagroup.com/en-US/faqs-corporate-information}{Qwen / 阿里巴巴 · 中国 36}}
\Box{108}{231.6}{87}{\fontsize{8.7}{12}\selectfont\color{muted}\sffamily \href{https://www.baai.ac.cn/about-us}{智源 BAAI · 中国 1}}
\Box{108}{238.2}{87}{\fontsize{8.7}{12}\selectfont\color{muted}\sffamily \href{https://bfl.ai/legal/imprint}{Black Forest Labs · 德国 1}}
\Box{108}{244.8}{87}{\fontsize{8.7}{12}\selectfont\color{muted}\sffamily \href{https://find-and-update.company-information.service.gov.uk/company/12295325}{Stability AI · 英国 1}}
\Box{108}{251.4}{87}{\fontsize{8.7}{12}\selectfont\color{muted}\sffamily \href{https://mistral.ai/fr/legal/}{Mistral AI · 法国 2}}
\Box{108}{258.0}{87}{\fontsize{8.7}{12}\selectfont\color{muted}\sffamily \href{https://www.onomaai.com/}{Onoma AI · 韩国 1}}
\Box{108}{264.6}{87}{\fontsize{8.7}{12}\selectfont\color{muted}\sffamily \href{https://www.lightricks.com/contact-us}{Lightricks · 以色列 1}}
\Box{108}{271.2}{87}{\fontsize{8.7}{12}\selectfont\color{muted}\sffamily \href{https://www.tii.ae/contact-us}{TII · 阿联酋 1}}
\end{canvas}
\begin{canvas}\Header{附录 A · 第三章方法与证据}{violet}\Footer{A4}
\Box{15}{27}{180}{\fontsize{20}{27}\selectfont\color{ink}\sffamily \bfseries 附录 A　案例与外部研究来源}
\Box{15}{43}{180}{\fontsize{9.5}{14.5}\selectfont\color{muted}\sffamily 机构公告证明组织安排，模型卡说明公开的来源关系，研究报告提供限定样本中的观察。资料的用途与时间分别列明，不据组织目标或资源数量推断实际公共成效。}
\Box{15}{63}{180}{\fontsize{9.4}{13.8}\selectfont\color{ink}\sffamily \bfseries \href{https://www.linuxfoundation.org/press/linux-foundation-announces-the-formation-of-the-agentic-ai-foundation}{LF：AAIF 成立公告}}
\Box{15}{69}{180}{\fontsize{8.8}{12.8}\selectfont\color{muted}\sffamily 2025-12-09；三个创始项目与组织定位。图 3-10。}
\Box{15}{79}{180}{\fontsize{9.4}{13.8}\selectfont\color{ink}\sffamily \bfseries \href{https://aaif.io/blog/agentgateway-joins-aaif-as-an-open-gateway-for-agentic-ai-infrastructure}{AAIF：agentgateway 加入公告}}
\Box{15}{85}{180}{\fontsize{8.8}{12.8}\selectfont\color{muted}\sffamily 2026-06-04；主期末扩至四项目。图 3-10。}
\Box{15}{95}{180}{\sffamily\bfseries\fontsize{9.4}{13.8}\selectfont \href{https://aaif.io/blog/a2a-joins-aaif}{AAIF：A2A 加入}\href{https://aaif.io/blog/agent-router-joins-aaif}{、Agent Router 加入}\href{https://aaif.io/news/agentic-ai-foundation-announces-four-new-projects}{及四项新增项目}}
\Box{15}{101}{180}{\fontsize{8.8}{12.8}\selectfont\color{muted}\sffamily 分别为 2026 年 8、9、10 月；期后扩展，当前目录为十项目。}
\Box{15}{111}{180}{\fontsize{9.4}{13.8}\selectfont\color{ink}\sffamily \bfseries \href{https://github.com/aaif/technical-committee/blob/main/governance/project-lifecycle-policy.md}{AAIF：项目生命周期政策}}
\Box{15}{117}{180}{\fontsize{8.8}{12.8}\selectfont\color{muted}\sffamily 采用 2026-10-15 版；治理门槛不回溯为整个主期的制度。}
\Box{15}{127}{180}{\fontsize{9.4}{13.8}\selectfont\color{ink}\sffamily \bfseries \href{https://www.linuxfoundation.org/press/agentic-ai-foundation-welcomes-the-gates-foundation-to-help-expand-the-benefits-of-agentic-ai-globally}{LF：Gates Foundation 加入 AAIF}}
\Box{15}{133}{180}{\fontsize{8.8}{12.8}\selectfont\color{muted}\sffamily 2026-09-15；支持非洲开发者的安排与承诺，不作为成效证据。}
\Box{15}{143}{180}{\fontsize{9.4}{13.8}\selectfont\color{ink}\sffamily \bfseries \href{https://atomgit.com/cann}{CANN 官方协作入口}}
\Box{15}{149}{180}{\fontsize{8.8}{12.8}\selectfont\color{muted}\sffamily 2026-11-10 观察；功能分层、社区、学习及算力入口。图 3-17。}
\Box{15}{159}{180}{\sffamily\bfseries\fontsize{9.4}{13.8}\selectfont \href{https://www.huawei.com/cn/news/2025/8/ascend-summit-CANN-open-source}{华为：CANN 开放公告}\href{https://www.huawei.com/cn/news/2025/9/hc-shengten-opensource}{与技术指导委员会公告}}
\Box{15}{165}{180}{\fontsize{8.8}{12.8}\selectfont\color{muted}\sffamily 2025-08-05、2025-09-20；开放方向、治理与上游协作安排。}
\Box{15}{175}{180}{\fontsize{9.4}{13.8}\selectfont\color{ink}\sffamily \bfseries \href{https://arxiv.org/html/2604.07190v1}{Lambert \& Brand：The ATOM Report}}
\Box{15}{181}{180}{\fontsize{8.8}{12.8}\selectfont\color{muted}\sffamily 2026-04-08；约 1,500 个重点开放权重语言模型，观察窗至 2026 年 3 月。}
\Box{15}{191}{180}{\fontsize{9.4}{13.8}\selectfont\color{ink}\sffamily \bfseries \href{https://github.com/huggingface/blog/blob/9e982d5f2ff0c93248a897bf1e0c59c9828c5ab8/state-of-open-models-summer-2026.md}{Hugging Face：Summer 2026 Observations}}
\Box{15}{197}{180}{\fontsize{8.8}{12.8}\selectfont\color{muted}\sffamily 2026-08-17 固定版；Qwen 衍生与 GGUF 转换，含主期后的观察。}
\Box{15}{207}{180}{\sffamily\bfseries\fontsize{9.4}{13.8}\selectfont \href{https://huggingface.co/deepseek-ai/DeepSeek-R1}{DeepSeek-R1}\href{https://huggingface.co/deepseek-ai/DeepSeek-R1-0528-Qwen3-8B}{ 与 R1-0528-Qwen3-8B 模型卡}}
\Box{15}{213}{180}{\fontsize{8.8}{12.8}\selectfont\color{muted}\sffamily 六个蒸馏版本及后续 Qwen3 接续；分别识别基础权重和推理样本关系。}
\Box{15}{223}{180}{\fontsize{9.4}{13.8}\selectfont\color{ink}\sffamily \bfseries \href{https://huggingface.co/unsloth/Qwen2.5-7B-Instruct}{Unsloth / Qwen2.5-7B-Instruct 模型卡}}
\Box{15}{229}{180}{\fontsize{8.8}{12.8}\selectfont\color{muted}\sffamily 基础模型元数据与发布版本；图 3-20 的适配接续例证。}
\Box{15}{239}{180}{\fontsize{9.4}{13.8}\selectfont\color{ink}\sffamily \bfseries \href{https://huggingface.co/BAAI/bge-m3}{BAAI / bge-m3 模型卡}}
\Box{15}{245}{180}{\fontsize{8.8}{12.8}\selectfont\color{muted}\sffamily 多语言检索与文本表示用途；图 3-19 的科研供给例证。}
\Box{15}{260}{180}{\fontsize{9.2}{13.8}\selectfont\color{muted}\sffamily ATOM 的下载份额分别以模型规模组为分母；OpenRouter 观察限每月前十开放模型。下载、调用 token（文本处理单位）与派生仓库均不等于独立用户或公共成效。}
\end{canvas}
\end{document}